
\section{9. Discussion}

\subsection{9.1 Key Takeaways}

\subsubsection{9.1.1 Takeaway 1: Union of Local + Global}

NSA's union of a fixed window and dynamic top-$k$ captures both regimes. 
Neither alone matches the union's accuracy.

\subsubsection{9.1.2 Takeaway 2: Attention Density is Low}

Even at 63\% density (short sequences), NSA preserves 95\% of the attention 
mass. At long sequences, density drops below 10\%.

\subsubsection{9.1.3 Takeaway 3: Zero Parameters Added}

NSA adds no parameters. Its benefit is \textbf{attention sparsity}, realized as 
lower memory and (with the right kernels) lower time.

\subsection{9.2 When to Use NSA}

\subsubsection{Recommended}

\begin{itemize}
\item \textbf{Long-context inference} where sparse kernels are available
\item \textbf{Streaming / KV cache} deployment (natural fit)
\item \textbf{Interpretability studies} of head specialization
\item \textbf{Base for kernel development} (sparse matmul)
\end{itemize}

\subsubsection{Not Recommended}

\begin{itemize}
\item \textbf{Short sequences} where density is high
\item \textbf{Parameter-constrained} applications
\item \textbf{Environments without sparse kernels}
\end{itemize}

\subsection{9.3 Choosing Window Size}

\begin{table}[H]
\centering
\small
\begin{tabular}{ll}
\toprule
\textbf{Sequence length} & \textbf{Recommended $w$} \\
\midrule
128 & 16 \\
512 & 32 \\
2048 & 64 \\
8192 & 128 \\
\bottomrule
\end{tabular}
\end{table}

\subsection{9.4 Choosing Top-K}

\begin{table}[H]
\centering
\small
\begin{tabular}{ll}
\toprule
\textbf{Use case} & \textbf{Recommended $k$} \\
\midrule
Memory-efficient & 4 \\
Balanced & 8 \\
High quality & 16 \\
Aggressive capture & 32 \\
\bottomrule
\end{tabular}
\end{table}

\subsection{9.5 Cost-Benefit}

For DistilBERT at $L = 128$:

\begin{table}[H]
\centering
\small
\begin{tabular}{lll}
\toprule
\textbf{Config} & \textbf{Accuracy} & \textbf{Density} \\
\midrule
Baseline & 90.71\% & 100\% \\
NSA ($w=16, k=8$) & 88.5 & 25\% \\
NSA ($w=32, k=16$) & 89.5 & 63\% \\
NSA ($w=64, k=32$) & 90.0 & 80\% \\
\bottomrule
\end{tabular}
\end{table}

\subsection{9.6 Interaction with Other Methods}

\begin{table}[H]
\centering
\small
\begin{tabular}{ll}
\toprule
\textbf{Method} & \textbf{NSA-compatible?} \\
\midrule
GQA & (multiplicative KV reduction) \\
RoPE & (RoPE before NSA) \\
Quantization & (INT8 NSA) \\
KV cache & (NSA reduces cache size) \\
MoE & (orthogonal) \\
\bottomrule
\end{tabular}
\end{table}

\subsection{9.7 Open Questions}

1. \textbf{Learned window and top-k:} Can we learn $w$ and $k$ per head?
2. \textbf{Sparse kernels:} When will they be widely available?
3. \textbf{Cross-layer selection transfer:} Do top-$k$ sets overlap across layers?
4. \textbf{Long-context evaluation:} How does NSA scale at $L > 8192$?

\subsection{9.8 Broader Impact}

\subsubsection{Positive}
\begin{itemize}
\item \textbf{Long-context inference:} Enables more capable applications.
\item \textbf{Cheaper inference:} Reduces compute per token (with sparse kernels).
\item \textbf{Interpretability:} Selection analysis reveals structure.
\end{itemize}

\subsubsection{Negative}
\begin{itemize}
\item \textbf{Dual-use:} Efficiency gains benefit malicious applications.
\item \textbf{Concentration:} Favors large labs.
\end{itemize}

\subsubsection{Environmental}
Total compute: ~10 GPU-hours, ~{{CO2}} kg CO$_2$eq.